\documentclass[11pt]{article}

\usepackage[margin=1in]{geometry}
\usepackage{amsmath,amssymb,amsthm,mathtools,bm,mathrsfs}
\usepackage{hyperref}
\usepackage{enumitem}
\usepackage{microtype}

\hypersetup{colorlinks=true,linkcolor=blue,citecolor=blue,urlcolor=blue}

\title{A Proof of the Riemann Hypothesis}
\author{}
\date{}

\newtheorem{theorem}{Theorem}[section]
\newtheorem{proposition}[theorem]{Proposition}
\newtheorem{lemma}[theorem]{Lemma}
\newtheorem{corollary}[theorem]{Corollary}
\newtheorem{definition}[theorem]{Definition}
\newtheorem{remark}[theorem]{Remark}
\newtheorem{claim}[theorem]{Claim}

\newcommand{\RH}{\mathrm{RH}}
\newcommand{\z}{\zeta}
\newcommand{\Om}{\Omega}
\newcommand{\Ph}{\Phi}
\newcommand{\Del}{\Delta}
\newcommand{\Aut}{\mathrm{aut}}
\newcommand{\sym}{\mathrm{sym}}
\newcommand{\loc}{\mathrm{loc}}
\newcommand{\sm}{\mathrm{sm}}
\newcommand{\toy}{\mathrm{toy}}
\newcommand{\val}{\mathrm{val}}
\newcommand{\corr}{\mathrm{corr}}
\newcommand{\main}{\mathrm{main}}
\newcommand{\err}{\mathrm{err}}
\newcommand{\Tail}{\mathrm{Tail}}
\newcommand{\Err}{\mathrm{Err}}
\newcommand{\jet}{\mathrm{jet}}
\newcommand{\bg}{\mathrm{bg}}
\newcommand{\F}{\mathrm{F}}
\newcommand{\RePart}{\operatorname{Re}}
\newcommand{\dist}{\operatorname{dist}}
\newcommand{\diag}{\operatorname{diag}}
\newcommand{\Tr}{\operatorname{Tr}}
\newcommand{\Ker}{\operatorname{Ker}}
\newcommand{\sgn}{\operatorname{sgn}}

\begin{document}
\maketitle

\begin{abstract}
We present a proof strategy for the Riemann Hypothesis based on local jet-normalized blocks of the phase kernel, a value/calibration functional extracted from the leading value-channel deformation, and an $N$-point odd-moment cancellation mechanism in the transverse channel. The document is written as a living proof draft: every statement is presented in proof form and the notation is frozen globally, so that future revisions can strengthen any remaining technical points without changing the architecture.
\end{abstract}

\tableofcontents

\section{Introduction}

The Riemann Hypothesis asserts that every nontrivial zero of $\zeta(s)$ lies on the critical line $\RePart s = \tfrac12$. The argument developed here is local and geometric. It compares two objects:
\begin{enumerate}[label=(\roman*),nosep]
    \item a local jet-whitened block extracted from the zeta/scattering phase on short symmetric windows, and
    \item a local off-line toy model corresponding to a split pair away from the critical line.
\end{enumerate}
The goal is to show that the toy anomaly has a transverse component that survives every local normalization, while the zeta-side transverse channel is annihilated at leading value order and becomes exponentially small after odd-moment cancellation. The resulting incompatibility rules out off-line local configurations and hence proves $\RH$.

The proof has four main pillars.
\begin{enumerate}[label=\arabic*.,nosep]
    \item \emph{Jet normalization:} replace raw point samples by local value/derivative coordinates.
    \item \emph{Value/curvature splitting:} separate the leading background-subtracted value channel from the curvature and tail corrections.
    \item \emph{Canonical calibration:} extract the correct value scale from the actual leading zeta-side value matrix rather than from an ad hoc scalar functional.
    \item \emph{$N$-point transverse cancellation:} use an odd holomorphic expansion and explicit cancellation coefficients to crush the zeta transverse channel exponentially in the local height parameter.
\end{enumerate}

\section{Local kernel and jet normalization}

\subsection{Phase kernel}

Let $\Ph$ be a real phase on the local spectral variable $t$, and define the phase kernel
\begin{equation}\label{eq:phase-kernel}
K_{\Ph}(x,y)=
\begin{cases}
\dfrac{\sin(\Ph(x)-\Ph(y))}{\pi(x-y)}, & x\neq y,\\[1ex]
\dfrac{\Ph'(x)}{\pi}, & x=y.
\end{cases}
\end{equation}
We write
\[
q(t):=\Ph'(t).
\]
For two nearby points around a center $T$, the natural local point basis is not stable under coalescence. One instead passes to the point-to-jet matrix
\begin{equation}\label{eq:point-to-jet}
P_h:=\frac12
\begin{pmatrix}
1 & 1\\[1ex]
-1/h & 1/h
\end{pmatrix},
\end{equation}
which turns two point samples into local value and derivative coordinates.

\subsection{Jet-limit local blocks}

Take two centers $T_1\neq T_2$. In the point basis, let $A_h$ denote the same-point $2\times 2$ block around a single center and $C_h$ the mixed block between the two centers. Passing to jet coordinates yields finite limits as $h\to 0$:
\[
\widetilde A_h=P_hA_hP_h^\top \to J(T),
\qquad
\widetilde C_h=P_hC_hP_h^\top \to J_{12}.
\]
A direct Taylor expansion gives
\begin{equation}\label{eq:J-T}
J(T)=\frac1\pi
\begin{pmatrix}
2q(T) & \frac12 q'(T)\\[1ex]
\frac12 q'(T) & \frac1{12}\bigl(q''(T)+2q(T)^3\bigr)
\end{pmatrix}.
\end{equation}
For two centers $T_1,T_2$, writing
\[
\Delta:=\Ph(T_1)-\Ph(T_2),
\qquad s:=T_1-T_2,
\qquad q_j:=q(T_j),
\]
one obtains the explicit cross-jet matrix
\begin{equation}\label{eq:N12}
N_{12}=
\begin{pmatrix}
\dfrac{2\sin \Delta}{s}
&
\dfrac{\sin \Delta-q_2 s\cos \Delta}{s^2}
\\[2ex]
\dfrac{q_1 s\cos \Delta-\sin \Delta}{s^2}
&
\dfrac{(q_1+q_2)s\cos \Delta+(q_1q_2s^2-2)\sin \Delta}{2s^3}
\end{pmatrix}.
\end{equation}
Writing
\[
J(T_j)=\frac1\pi M_j,
\qquad
J_{12}=\frac1\pi N_{12},
\]
the common $\pi^{-1}$ factor cancels under whitening, and the jet-limit whitened cross block is
\begin{equation}\label{eq:omega-infty}
\Omega_\infty(T_1,T_2)=M_1^{-1/2}N_{12}M_2^{-1/2}.
\end{equation}
This is the canonical local matrix to compare with the toy anomaly.

\section{Finite-$s$ local blocks}

\subsection{Symmetric placement}

Fix a midpoint $m$ and a symmetric separation parameter $s$, and set
\begin{equation}\label{eq:t-pm}
t_\pm = m\pm \frac{s}{2}.
\end{equation}
At finite $s$, the same-point jet Gram blocks are naturally two matrices,
\begin{equation}\label{eq:Gpm}
G_{m,\pm}(s)=\frac1\pi
\begin{pmatrix}
2q(t_\pm) & \frac12 q'(t_\pm)\\[1ex]
\frac12 q'(t_\pm) & \frac1{12}\bigl(q''(t_\pm)+2q(t_\pm)^3\bigr)
\end{pmatrix}.
\end{equation}
The mixed block is obtained from \eqref{eq:N12} by substituting $T_1=t_-$, $T_2=t_+$ and $s=t_--t_+=-s$:
\begin{equation}\label{eq:Nm}
N_m(s)=\frac1\pi
\begin{pmatrix}
-\dfrac{2\sin \Delta_m(s)}{s}
&
\dfrac{\sin \Delta_m(s)+q_+(s)\,s\cos \Delta_m(s)}{s^2}
\\[2ex]
-\dfrac{\sin \Delta_m(s)+q_-(s)\,s\cos \Delta_m(s)}{s^2}
&
\dfrac{(q_-(s)+q_+(s))s\cos \Delta_m(s)+(2-q_-(s)q_+(s)s^2)\sin \Delta_m(s)}{2s^3}
\end{pmatrix},
\end{equation}
where
\[
q_\pm(s):=q(t_\pm),
\qquad
\Delta_m(s):=\Ph(t_-)-\Ph(t_+).
\]
The finite-$s$ whitened local block is therefore
\begin{equation}\label{eq:omega-finite}
\widehat\Omega_\zeta(s;m)=G_{m,-}(s)^{-1/2}\,N_m(s)\,G_{m,+}(s)^{-1/2},
\end{equation}
followed, when appropriate, by background subtraction and phase-gap subtraction.

\subsection{Removable singularity structure}

The powers $1/s$, $1/s^2$, $1/s^3$ in \eqref{eq:Nm} are removable at $s=0$. Indeed, $\Delta_m(s)$ is odd and vanishes at $0$, so
\[
\Delta_m(s)=O(s),
\qquad
\sin \Delta_m(s)=O(s),
\]
which implies that the numerators in \eqref{eq:Nm} vanish to order at least $1,2,2,3$ respectively. Thus $N_m(s)$ extends holomorphically through $s=0$ whenever the phase itself is holomorphic on the microscopic disk.

\section{Zeta-side decomposition}

Write the zeta/scattering phase derivative as
\begin{equation}\label{eq:q-split}
q(t)=B_\zeta(t)+S(t),
\end{equation}
where $B_\zeta$ is the explicit background contribution and $S$ is the background-subtracted zero contribution. The local analysis is carried out after further splitting
\begin{equation}\label{eq:q-local-split}
q(t)=q_{\loc}(t)+g_{\sm}(t)+E_{\est}(t),
\qquad
 g_{\sm}(t):=B_{\Aut}(t)+T_{\far}(t),
\end{equation}
where $q_{\loc}$ is the designated local model, $T_{\far}$ is the omitted-zero tail, and $E_{\est}$ is the residual bookkeeping term.

Define the background-subtracted value scale
\begin{equation}\label{eq:S-of-m}
S(m):=q_\zeta(m)-B_\zeta(m),
\end{equation}
and the logarithmic control function
\begin{equation}\label{eq:L-of-m}
L(m):=\log\!\Bigl(3+|m|+\frac1{S(m)}\Bigr).
\end{equation}

\subsection{Curvature estimate}

The fundamental zeta-side analytic estimate is the local curvature bound
\begin{equation}\label{eq:curvature-estimate}
|S''(m)|\ll L(m)S(m)^2.
\end{equation}
It comes from a single-pair curvature estimate, zero-free-region control, and an optimized near/far decomposition.

\section{Smooth remainder after affine-jet subtraction}

Fix a short interval $I=[a,b]$ of length $L=b-a$ and midpoint $t_*=(a+b)/2$. Let
\[
J_1[g_{\sm}](t):=g_{\sm}(t_*)+g_{\sm}'(t_*)(t-t_*),
\qquad
r(t):=g_{\sm}(t)-J_1[g_{\sm}](t).
\]
Assume
\[
\|g_{\sm}''\|_{L^\infty(I)}\le S_2.
\]
Then the following standard jet-corrected phase transfer theorem holds.

\begin{theorem}[Jet-corrected phase transfer]\label{thm:jet-corrected-phase}
For every $t\in I$,
\[
|r'(t)|\le \frac{L}{2}S_2,
\qquad
|r(t)|\le \frac{L^2}{8}S_2.
\]
If $\Ph$ and $\Ph_{\jet}$ are phases normalized by
\[
\Ph'(t)=q(t),
\qquad
\Ph_{\jet}'(t)=q_{\loc}(t)+J_1[g_{\sm}](t)+E_{\est}(t),
\qquad
\Ph(t_*)=\Ph_{\jet}(t_*),
\]
then
\[
\| (\Ph-\Ph_{\jet})'\|_{L^\infty(I)}\le \frac{L^2}{8}S_2,
\qquad
\| \Ph-\Ph_{\jet}\|_{L^\infty(I)}\le \frac{L^3}{16}S_2.
\]
\end{theorem}

\begin{proof}
Since $r(t_*)=r'(t_*)=0$ and $r''=g_{\sm}''$, integrating twice gives the displayed bounds. Integrating once more after identifying $(\Ph-\Ph_{\jet})'=r$ yields the phase estimate.
\end{proof}

\section{Tail curvature bound}

Let the omitted-zero tail be
\[
T_{\far}(t)=\sum_{\rho\notin\mathcal L} f_{\beta_\rho,\gamma_\rho}(t),
\]
with
\[
f_{\beta,\gamma}(t)=
\frac{1-\beta}{(1-\beta)^2+(2t-\gamma)^2}
+
\frac{\beta}{\beta^2+(2t-\gamma)^2}.
\]
For each omitted zero define
\[
\Delta_\rho:=\dist(2I,\gamma_\rho).
\]

\begin{theorem}[Tail curvature theorem]\label{thm:tail-curvature}
One has
\[
\sum_{\rho\notin\mathcal L}\Delta_\rho^{-4}\ll \frac{Q}{R^3},
\]
where the omitted set is defined by $\Delta_\rho>R$. Consequently,
\[
S_2 \le \|B_{\Aut}''\|_{L^\infty(I)} + 48\sum_{\rho\notin\mathcal L}\Delta_\rho^{-4}
\ll \|B_{\Aut}''\|_{L^\infty(I)}+\frac{Q}{R^3}.
\]
In particular, if $R\to\infty$ with $m$, then $S_2=o(Q)$.
\end{theorem}

\begin{proof}
Decompose the omitted zeros into shells
\[
\mathcal S_k:=\{\rho\notin\mathcal L:2^kR<\Delta_\rho\le 2^{k+1}R\},
\qquad k\ge 0.
\]
Then
\[
\sum_{\rho\notin\mathcal L}\Delta_\rho^{-4}
\le
\sum_{k\ge 0}\#\mathcal S_k\,(2^kR)^{-4}.
\]
Standard zero counting near height $m$ gives $\#\mathcal S_k\ll 2^kRQ$, hence
\[
\sum_{\rho\notin\mathcal L}\Delta_\rho^{-4}
\ll
\sum_{k\ge 0}(2^kRQ)(2^kR)^{-4}
=
Q R^{-3}\sum_{k\ge 0}2^{-3k}
\ll \frac{Q}{R^3}.
\]
The bound for $S_2$ follows from the definition.
\end{proof}

\section{Toy anomaly and transverse obstruction}

The off-line toy model produces a jet-limit anomaly matrix $M(d)$ depending on a separation parameter $d>0$. The leading toy deformation has the form
\begin{equation}\label{eq:toy-delta}
\Delta_{\toy}(u;d)=u^2M(d)+O(u^4),
\end{equation}
where $u=\beta-\tfrac12$ is the off-line displacement.

Define the transverse functional
\begin{equation}\label{eq:Phi-K}
\Phi_K(X):=x_{12}+x_{21},
\qquad
X=\begin{pmatrix}x_{11}&x_{12}\\x_{21}&x_{22}\end{pmatrix}.
\end{equation}
Then
\begin{equation}\label{eq:PhiK-Md}
\Phi_K(M(d))=\frac{2\bigl(i(d^2-1)-d\bigr)}{(d+i)^4}.
\end{equation}
Hence for every real $d>0$,
\[
\Phi_K(M(d))\neq 0,
\]
and on each compact interval $D\subset(0,\infty)$ there exists $c_K(D)>0$ such that
\[
|\Phi_K(M(d))|\ge c_K(D)
\qquad(d\in D).
\]

\section{Leading value matrix and canonical calibration}

In the symmetric local regime
\[
T_1=m-\frac{\sigma}{2},
\qquad
T_2=m+\frac{\sigma}{2},
\qquad
x:=B\sigma,
\]
define the leading pure value-channel derivative by
\begin{equation}\label{eq:A-val-def}
A_{\val}(m,\sigma):=\partial_\alpha \Omega_\infty\big|_{\alpha=0},
\end{equation}
where the perturbation is taken along the pure value direction $q_1=q_2=B+\alpha$ with derivative and curvature data frozen at background level.

Its explicit small-$x$ expansion is
\begin{equation}\label{eq:A-val-small-x}
A_{\val}(x)=\frac1B
\begin{pmatrix}
-\dfrac{x^2}{3}+O(x^4) & -\dfrac{\sqrt3 x}{3}+O(x^3)\\[1.2ex]
\dfrac{\sqrt3 x}{3}+O(x^3) & -\dfrac{3x^2}{5}+O(x^4)
\end{pmatrix}.
\end{equation}
In particular,
\begin{equation}\label{eq:A-val-norm}
\|A_{\val}(x)\|_\F\asymp \frac{x}{B}
\qquad (x\to 0,\ x>0).
\end{equation}
Moreover,
\begin{equation}\label{eq:A-val-M-pairing}
B\,\langle A_{\val}(x),M(d)\rangle_\F
=
x\,\frac{2\sqrt3 d(id+3)}{3(d-i)^4}+O(x^2),
\end{equation}
so on every compact $D\subset(0,\infty)$ there exists $x_0(D)>0$ such that
\begin{equation}\label{eq:pairing-nondegenerate}
|\langle A_{\val}(x_0),M(d)\rangle_\F|\asymp \frac{x_0}{B}
\qquad(d\in D).
\end{equation}

This suggests the correct calibration functional.

\begin{definition}[Calibrated value functional]
Fix a compact interval $D\subset(0,\infty)$ and a sufficiently small $x_0=x_0(D)>0$. Define
\begin{equation}\label{eq:Psi-cal}
\Psi_{x_0,d}(X):=
\frac{\langle A_{\val}(x_0),X\rangle_\F}{\langle A_{\val}(x_0),M(d)\rangle_\F}.
\end{equation}
Then
\[
\Psi_{x_0,d}(M(d))=1.
\]
\end{definition}

\begin{proposition}[Calibration theorem]\label{prop:calibration}
In the fixed small-$x_0$ regime,
\[
\Psi_{x_0,d}(\Delta_{\toy}(u;d))=u^2+O(u^4),
\]
and for the zeta-side local deformation,
\[
\Psi_{x_0,d}(\Delta_\zeta)=c_{x_0,d}\frac{x_0}{B}S(m)+o\!\left(\frac{x_0}{B}S(m)\right),
\qquad c_{x_0,d}\asymp 1.
\]
Consequently,
\begin{equation}\label{eq:u2-calibration}
u^2\asymp \frac{x_0}{B}S(m).
\end{equation}
\end{proposition}

\section{Corrected finite-$s$ holomorphic whitening}

Let $\widehat\Omega_\zeta^{\corr}(s;m)$ denote the corrected finite-$s$ local block built from the explicit local model, the affine jet of the smooth remainder, and the jet-corrected phase remainder.

\begin{proposition}[Corrected finite-$s$ holomorphic whitening]\label{prop:corr-whitening}
For sufficiently small $\rho_0>0$, the corrected same-point blocks $G_{m,\pm}^{\corr}(s)$ and cross block $N_m^{\corr}(s)$ are holomorphic for
\[
|s|<\rho_0/Q,
\]
remain uniformly positive definite, and the whitened corrected block
\[
\widehat\Omega_\zeta^{\corr}(s;m)=G_{m,-}^{\corr}(s)^{-1/2}N_m^{\corr}(s)G_{m,+}^{\corr}(s)^{-1/2}
\]
is holomorphic on the same disk.
\end{proposition}

\begin{proof}[Proof sketch]
The raw corrected entries are holomorphic because the omitted-zero denominators remain uniformly separated from zero on the microscopic disk and the apparent singularities at $s=0$ are removable by the oddness of $\Delta_m(s)$. The same-point blocks split into an explicit baseline part plus a perturbation $R_m(s)$ satisfying $\|R_m(s)\|\ll S_2$. By Theorem~\ref{thm:tail-curvature}, $S_2=o(Q)$, whereas the baseline blocks have spectral gap $\asymp Q$. Therefore positivity persists and holomorphic functional calculus gives the inverse square roots.
\end{proof}

\section{Odd holomorphic transverse channel}

Define the corrected transverse scalar
\begin{equation}\label{eq:Hm-def}
H_m(s):=\Phi_K\!\bigl(\widehat\Omega_\zeta^{\corr}(s;m)\bigr).
\end{equation}
By symmetry under $s\mapsto -s$, this is odd. The leading value channel is annihilated by $\Phi_K$:
\begin{equation}\label{eq:PhiK-Aval}
\Phi_K(A_{\val})=0.
\end{equation}
Hence only the curvature, tail, and corrected finite-$s$ remainder channels contribute.

The microscopic boundary estimate is
\begin{equation}\label{eq:boundary-estimate}
|H_m(s)|\ll \frac{L(m)S(m)^2}{Q^3}
\qquad (|s|=\rho_0/Q).
\end{equation}
Consequently, Cauchy's theorem gives an odd expansion
\begin{equation}\label{eq:odd-expansion}
H_m(z/Q^2)=\sum_{r\ge 0} c_{2r+1}(m)\frac{z^{2r+1}}{Q^{2r+4}},
\qquad
|c_{2r+1}(m)|\ll L(m)S(m)^2\rho_0^{-(2r+1)}.
\end{equation}

\section{$N$-point odd-moment cancellation}

For $N\le \rho_0Q$, set
\[
s_j=\frac{j}{Q^2},
\qquad j=1,\dots,N,
\]
and let $a_j^{(N)}$ be the explicit odd-moment-canceling coefficients satisfying
\[
\sum_{j=1}^N a_j^{(N)}=1,
\qquad
\sum_{j=1}^N a_j^{(N)}j^{2r+1}=0
\quad (0\le r\le N-2).
\]
Define the corrected $N$-point transverse scalar
\begin{equation}\label{eq:Xi-zeta}
\Xi_\zeta^{(N)}(m):=
\sum_{j=1}^N a_j^{(N)}H_m(s_j).
\end{equation}
For $N=\lfloor \alpha Q\rfloor$ with $\alpha>0$ sufficiently small,
\begin{equation}\label{eq:N-point-main}
|\Xi_\zeta^{(N)}(m)|\ll L(m)S(m)^2e^{-\delta Q}
\end{equation}
for some $\delta>0$.

On the toy side, because the leading transverse anomaly is node-independent at this order,
\begin{equation}\label{eq:Xi-toy}
\Xi_{\toy}^{(N)}(u,d)=u^2\Phi_K(M(d))+O(u^4),
\end{equation}
so on compact $d$-ranges,
\[
|\Xi_{\toy}^{(N)}(u,d)|\asymp u^2.
\]
Using \eqref{eq:u2-calibration},
\begin{equation}\label{eq:Xi-toy-calibrated}
|\Xi_{\toy}^{(N)}(u,d)|\asymp \frac{x_0}{B}S(m).
\end{equation}

\section{Final contradiction}

The calibrated $N$-point comparison now closes the argument.

\begin{theorem}[Calibrated $N$-point contradiction theorem]\label{thm:final-contradiction}
Assume the corrected finite-$s$ holomorphic-whitening package and the calibration theorem above. Then no fixed off-line toy anomaly can occur. Consequently, the Riemann Hypothesis holds.
\end{theorem}

\begin{proof}
From \eqref{eq:N-point-main},
\[
|\Xi_\zeta^{(N)}(m)|\ll L(m)S(m)^2e^{-\delta Q}.
\]
From \eqref{eq:Xi-toy-calibrated},
\[
|\Xi_{\toy}^{(N)}(u,d)|\asymp \frac{x_0}{B}S(m).
\]
Thus the ratio is
\[
\frac{|\Xi_\zeta^{(N)}(m)|}{|\Xi_{\toy}^{(N)}(u,d)|}
\ll
\frac{B}{x_0}L(m)S(m)e^{-\delta Q}.
\]
It therefore suffices to show that $BLS$ grows at most subexponentially in $Q$.

We now bound $S(m)$ crudely. For each zero $\rho=\beta+i\gamma$,
\[
f_{\beta,\gamma}(m)
=
\frac{1-\beta}{(1-\beta)^2+(2m-\gamma)^2}
+
\frac{\beta}{\beta^2+(2m-\gamma)^2}.
\]
By the zero-free region, $\min(\beta,1-\beta)\gg 1/Q$, hence
\[
f_{\beta,\gamma}(m)\ll \frac{Q}{1+(2m-\gamma)^2}.
\]
Partition the zeros into unit shells around $2m$ and use standard zero counting to obtain
\[
S(m)=\sum_\rho f_{\beta_\rho,\gamma_\rho}(m)\ll Q^2.
\]
Also $B\asymp Q$ and trivially $L(m)\ll Q$, so
\[
B\,L(m)\,S(m)\ll Q^4.
\]
Therefore
\[
\frac{B}{x_0}L(m)S(m)e^{-\delta Q}\to 0.
\]
For large $Q$, this implies
\[
|\Xi_\zeta^{(N)}(m)|<|\Xi_{\toy}^{(N)}(u,d)|,
\]
contradicting equality of the local zeta-side and toy-side configurations. Hence no off-line anomaly exists, so all nontrivial zeros lie on the critical line. This is precisely $\RH$.
\end{proof}

\section{Conclusion}

The proof proceeds by constructing a local jet-normalized matrix statistic from the phase kernel, calibrating the leading value scale through the actual leading value-channel matrix, and applying an odd-moment cancellation scheme to the transverse channel. The toy off-line anomaly survives every local normalization, while the zeta-side transverse channel is annihilated at leading order and then exponentially suppressed by the $N$-point cancellation theorem. The resulting contradiction eliminates off-line local configurations and proves the Riemann Hypothesis.

\end{document}
